\documentclass[12pt,a4paper]{article}
\usepackage[spanish]{babel}
\usepackage[utf8]{inputenc}
\usepackage[T1]{fontenc}
\usepackage[a4paper,margin=2.5cm]{geometry}
\usepackage{listings}
\usepackage{xcolor}
\usepackage{parskip}

\title{Actividad evaluativa 1\\Centro de monitoreo de accesos}
\author{Nombre y apellidos: Carmen Maruri Pujadas\\}
\date{23 de octubre de 2026}

\begin{document}
\maketitle

\section{Objetivo}
El objetivo de la actividad es desarrollar en Java un programa que analice
dos servidores y determine su estado de seguridad a partir de los intentos
de autenticación. El programa valida los datos, calcula el porcentaje de
fallos cuando corresponde y compara la prioridad de los dos servidores.

\section{Diseño del programa}
El programa está organizado en dos clases. La clase \texttt{Servidor}
contiene los cinco atributos solicitados y los métodos necesarios para
validar los datos, evaluar el estado, calcular el porcentaje de fallos,
mostrar el resumen y obtener el nivel de prioridad. La clase \texttt{Main}
lee los datos con \texttt{Scanner}, crea dos objetos independientes y
presenta el informe final.

\section{Validación y clasificación}
Antes de clasificar un servidor se comprueba que el nombre no esté vacío,
que el origen sea interno o externo, que la criticidad sea si o no, que el
total de intentos sea mayor que cero y que los fallos estén entre cero y
el total. Las comparaciones de texto ignoran las mayúsculas y los espacios
al principio o al final.

Si los datos son válidos, se aplican las reglas en el orden indicado en el
enunciado: \texttt{BLOQUEADO}, \texttt{ALERTA} y \texttt{NORMAL}. Si los
datos no son válidos, se muestra \texttt{DATOS INVALIDOS} y el porcentaje
aparece como \texttt{N/D}. La comparación final no se realiza si alguno
de los servidores tiene datos inválidos.

\section{Pruebas}
Se deben ejecutar por separado los casos de prueba del enunciado. Los
resultados esperados son:

\begin{center}
\begin{tabular}{|l|l|l|r|r|l|}
\hline
Nombre & Origen & Crítico & Total & Fallos & Estado esperado \\
\hline
VPN & externo & si & 10 & 3 & BLOQUEADO \\
WEB & interno & no & 8 & 4 & ALERTA \\
CORREO & externo & no & 8 & 2 & ALERTA \\
VPN & interno & si & 12 & 2 & NORMAL \\
WEB & externo & no & 5 & 5 & BLOQUEADO \\
WEB & interno & no & 4 & 5 & DATOS INVALIDOS \\
WEB & interno & no & 0 & 0 & DATOS INVALIDOS \\
CORREO & local & no & 10 & 2 & DATOS INVALIDOS \\
WEB & externo & quizás & 10 & 2 & DATOS INVALIDOS \\
\hline
\end{tabular}
\end{center}

También hay que probar entradas con espacios al principio y al final,
por ejemplo \texttt{ EXTERNO } y \texttt{ SI }, para comprobar que se
aceptan correctamente.

\section{Conclusión}
La solución permite revisar dos servidores con las mismas reglas y
comparar su nivel de riesgo. Separar la validación, la clasificación y
la presentación en métodos facilita la lectura del código y evita
calcular porcentajes cuando los datos no son válidos.

\end{document}
